\section{Sistematika Penulisan}

Penulisan tesis ini disusun dalam lima bab sebagai berikut.

Bab I Pendahuluan membahas latar belakang, rumusan masalah, tujuan penelitian, hipotesis, batasan masalah, metodologi penelitian secara ringkas, dan sistematika penulisan.

Bab II Tinjauan Pustaka dan Landasan Teori membahas landasan teori mengenai \textit{Internet of Things} dan keterbatasan sumber daya perangkat \textit{edge}, kontrol akses berbasis atribut (ABAC) beserta penambangan kebijakannya, algoritma metaheuristik dan \textit{Estimation of Distribution Algorithm} (EDA) beserta variannya yang ringan dan sadar dependensi, posisi penelitian terhadap literatur terdahulu, argumentasi kebaruan, serta kerangka berpikir penelitian.

Bab III Metodologi Penelitian menjelaskan pendekatan \textit{Design Science Research} yang digunakan, tahapan penelitian, variabel penelitian, rancangan algoritma CoDA-EDA yang diusulkan, konstruksi dataset, algoritma pembanding, perangkat dan lingkunga eksperimen, desain eksperimen, metrik dan instrumen evaluasi, teknik analisis data, serta studi pendahuluan yang mencakup analisis karakteristik dataset dan kalibrasi rancangan.

Bab IV Hasil dan Pembahasan menyajikan hasil pengujian empiris keempat hipotesis pada kelima dataset ABAC Lab, validasi eksternal pada datset \textins{Smart Building IoT}, contoh kebijakan hasil \textit{mining}, serta sintesis temuan, keterbatasan, dan jawaban atas rumusan masalah.

Bab V Kesimpulan dan Saran memuat memuat kesimpulan penelitian yang menjawab tujuan dan rumusan masalah beserta kontribusi yang dicapai, serta saran penerapan praktis dan pengembangan lanjutan.